\documentclass[10pt]{beamer}

% ---------------- THEME ----------------
\usetheme{Madrid}
\usecolortheme{seahorse}
\setbeamertemplate{navigation symbols}{}

% ---------------- PACKAGES ----------------
\usepackage{amsmath, amssymb}
\usepackage{graphicx}
\usepackage{booktabs}

% ---------------- TITLE ----------------
\title[DRL Trading Framework]{Deep Reinforcement Learning for Multi-Asset Trading}
\author{Rousomanis Georgios (10703)}
\date{May 2026}

\begin{document}

% =====================================================
\begin{frame}
\titlepage
\end{frame}

% =====================================================
\begin{frame}{Motivation}
\begin{itemize}
    \item Financial markets are sequential, stochastic, and non-stationary
    \item Classical strategies struggle under regime shifts
    \item DRL enables adaptive decision-making directly from data
\end{itemize}
\end{frame}

% =====================================================
\begin{frame}{System Overview}
\begin{block}{Full Pipeline}
\begin{itemize}
    \item Custom multi-asset trading environments
    \item Deep Reinforcement Learning agents
    \item Temporal feature modeling (MLP, LSTM, Transformer)
    \item Risk control (Fuzzy Logic)
    \item Hyperparameter optimization (Genetic Algorithm)
\end{itemize}
\end{block}
\end{frame}

% =====================================================
\begin{frame}{Trading Environment}
\begin{itemize}
    \item Multi-asset portfolio (stocks + crypto)
    \item Gymnasium-based sequential decision process
    \item Observation contains:
    \begin{itemize}
        \item Portfolio state (cash + holdings)
        \item Market state (prices + indicators)
        \item Risk signals (VIX, turbulence for stocks)
    \end{itemize}
\end{itemize}

\pause

State representation:
\[
s_t = [s_t^{portfolio}, s_t^{market}]
\]

\pause

Continuous action space:
\[
a_t \in [-1,1]^N
\]

\pause
Reward:
\[
r_t = \log\left(\frac{V_{t+1} + \epsilon}{V_t + \epsilon}\right)
 - \lambda e_t v_t
\]
\begin{itemize}
    \item $e_t$: portfolio exposure
    \item $v_t$: market volatility
    \item $\lambda$: risk sensitivity coefficient
\end{itemize}
\end{frame}

% =====================================================
\begin{frame}{Policy Architectures}
\begin{itemize}
    \item MLP (baseline)
    \item LSTM (temporal memory)
    \item Transformer (attention-based modeling)
    \item Hybrid MLP--Transformer
\end{itemize}

\begin{block}{Key Challenge}
Financial data requires modeling both:
\begin{itemize}
    \item temporal dependencies
    \item cross-asset interactions
\end{itemize}
\end{block}
\end{frame}

% =====================================================
\begin{frame}{Transformer Feature Extractor}

\textbf{Input Construction}

\[
X =
[x_{t-T+1}, \dots, x_t]
\]

where each observation contains both portfolio and market information:

\[
x_t =
[s_t^{portfolio}, s_t^{market}]
\]

\textbf{Architecture}

\begin{itemize}
    \item Sequence length: $T = 10$
    \item Linear projection to $d_{model}=128$
    \item Learnable positional embeddings
    \item Transformer Encoder:
    \begin{itemize}
        \item 2 encoder layers
        \item 4 attention heads
        \item Feed-forward dimension 256
        \item Dropout 0.1
    \end{itemize}
\end{itemize}

Final representation:

\[
h_t = Transformer(X)_T
\]

The embedding $h_t$ is provided to the actor and critic networks.

\end{frame}

% =====================================================
\begin{frame}{Hybrid MLP--Transformer Architecture}

\textbf{Motivation}

Market observations exhibit temporal structure, while portfolio information
represents an instantaneous allocation state.

\pause
\vspace{0.3cm}

\textbf{Market Branch}

\[
X^{market}
=
[s^{market}_{t-T+1}, \dots, s^{market}_{t}]
\]

\[
h_t = Transformer(X^{market})_T
\]

\pause
\vspace{0.2cm}

\textbf{Portfolio Branch}

\[
h_t^{portfolio}
=
MLP(s_t^{portfolio})
\]

\vspace{0.2cm}

\textbf{Fusion}

\[
z_t =
[h_t^{market}, h_t^{portfolio}]
\]

Concatenation followed by a projection layer before being passed to the actor
and critic networks.

\end{frame}

% =====================================================
\begin{frame}{Data Processing}
\begin{itemize}
    \item Rolling window sequences (T = 10)
    \item Online feature-wise normalization
    \item Reward normalization via rolling statistics
\end{itemize}

\begin{block}{Why it matters}
\begin{itemize}
    \item Handles non-stationary distributions
    \item Prevents scale dominance
\end{itemize}
\end{block}
\end{frame}

% =====================================================
\begin{frame}{Reinforcement Learning Algorithms}

\begin{columns}[T]

\column{0.48\textwidth}

\textbf{On-Policy Methods}

\vspace{0.2cm}

\textbf{A2C}
\begin{itemize}
    \item Actor-Critic architecture
    \item Uses advantage estimates to reduce variance
    \item Stable but less sample efficient
\end{itemize}

\vspace{0.2cm}

\textbf{PPO}
\begin{itemize}
    \item Clipped policy updates
    \item Prevents large policy shifts
    \item Widely used due to strong stability
\end{itemize}

\column{0.48\textwidth}

\textbf{Off-Policy Methods}

\vspace{0.2cm}

\textbf{DDPG}
\begin{itemize}
    \item Deterministic actor-critic
    \item Replay buffer for sample reuse
    \item Designed for continuous actions
\end{itemize}

\vspace{0.1cm}

\textbf{TD3}
\begin{itemize}
    \item Extension of DDPG
    \item Twin critics reduce overestimation bias
    \item Delayed policy updates improve stability
\end{itemize}

\vspace{0.1cm}

\textbf{SAC}
\begin{itemize}
    \item Stochastic actor-Critic
    \item Strong sample efficiency
\end{itemize}

\end{columns}

\end{frame}

% =====================================================
\begin{frame}{Training Performance}
\begin{figure}
    \centering
    \includegraphics[width=0.9\linewidth]{figures/rollout_ep_rew_mean.png}
    \caption{Training reward evolution}
\end{figure}
\end{frame}

% =====================================================
\begin{frame}{Financial Performance}
\begin{figure}
    \centering
    \includegraphics[width=0.9\linewidth]{figures/financial_total_return.png}
    \caption{Total return across models}
\end{figure}
\end{frame}

% =====================================================
\begin{frame}{Risk-Adjusted Performance}
\begin{figure}
    \centering
    \includegraphics[width=0.9\linewidth]{figures/financial_sharpe_ratio.png}
    \caption{Sharpe ratio comparison}
\end{figure}
\end{frame}

% =====================================================
\begin{frame}{Key Findings}
\begin{itemize}
    \item PPO $\rightarrow$ most stable learning behavior
    \item DDPG (MLP) $\rightarrow$ strongest overall performance
    \item Transformers $\rightarrow$ higher representation power, lower stability (off-policy)
    \item LSTM $\rightarrow$ strong but unstable long-horizon behavior
\end{itemize}
\end{frame}

% =====================================================
\begin{frame}{Backtesting Evaluation}

\begin{columns}

\column{0.4\textwidth}

\textbf{Evaluation Setup}

\vspace{0.2cm}

\begin{itemize}
    \item Out-of-sample period:
    \begin{itemize}
        \item Jan 2026 -- Mar 2026
    \end{itemize}

    \item Dow Jones 30 constituents

    \item Metrics:
    \begin{itemize}
        \item Total Return
        \item Portfolio Value
        \item Sharpe Ratio
    \end{itemize}
\end{itemize}

\vspace{0.3cm}

\textbf{Selected Models}
\begin{itemize}
    \item DDPG (MLP)
    \item PPO (Transformer)
    \item PPO (Hybrid)
\end{itemize}

\column{0.6\textwidth}

\begin{figure}
    \centering
    \includegraphics[width=\linewidth]{figures/backtest_total_return.png}
    \caption{Out-of-sample total return.}
\end{figure}

\end{columns}

\end{frame}

% =====================================================
\begin{frame}{Portfolio Value \& Risk-Adjusted Performance}

\begin{columns}

\column{0.5\textwidth}
\begin{figure}
    \centering
    \includegraphics[width=\linewidth]{figures/backtest_portfolio_value.png}
    \caption{Portfolio value evolution}
\end{figure}

\column{0.5\textwidth}
\begin{figure}
    \centering
    \includegraphics[width=\linewidth]{figures/backtest_sharpe_ratio.png}
    \caption{Sharpe ratio evolution}
\end{figure}

\end{columns}

\vspace{0.2cm}

\begin{block}{Key Observation}
DDPG (MLP) maintains the strongest portfolio growth and the most
consistent risk-adjusted performance during the out-of-sample period,
while both Transformer-based PPO variants exhibit gradual degradation.
\end{block}

\end{frame}

% =====================================================
\begin{frame}{Fuzzy-Enhanced Risk Management}

\textbf{Motivation}

\vspace{0.2cm}

The original environment relied on hard threshold-based controls:

\[
a_t' = \min(a_t,0)
\]

which completely blocked buying actions during periods of high market
turbulence.

\vspace{0.3cm}

\textbf{Limitation}

\begin{itemize}
    \item Abrupt action suppression
    \item Discontinuous behavior
    \item Difficult learning dynamics
\end{itemize}

\vspace{0.3cm}

\textbf{Proposed Solution}

A Mamdani Fuzzy Inference System composed of:

\begin{enumerate}
    \item Fuzzy Action Smoother
    \item Fuzzy Reward Shaper
\end{enumerate}

Both components replace hard decisions with smooth, context-dependent
risk adjustments.

\end{frame}

% =====================================================
\begin{frame}{Fuzzy Action Smoother}

\textbf{Goal}

Replace the hard turbulence constraint

\[
a_t'=\min(a_t,0)
\]

with a smooth risk-aware action adjustment mechanism.

\vspace{0.3cm}

\textbf{Inputs}

\begin{itemize}
    \item Raw Action ($a_i$):
    \[
    \{\text{Sell},\ \text{Hold},\ \text{Buy}\}
    \]

    \item Market Turbulence ($\tau_t$):
    \[
    \{\text{Low},\ \text{Medium},\ \text{High}\}
    \]
\end{itemize}

\vspace{0.3cm}

\textbf{Output}

\[
a_{t,\text{safe}}
\]

Risk-adjusted trading action obtained through Mamdani inference and
centroid defuzzification.

\end{frame}

% =====================================================
\begin{frame}{Fuzzy Action Smoother Rules}

\begin{columns}[T]

\column{0.42\textwidth}

\textbf{Rule Base}

\vspace{0.2cm}

\begin{enumerate}
    \item \textbf{IF} Turbulence is High\\
          \textbf{THEN} Action is Sell

    \vspace{0.2cm}

    \item \textbf{IF} Turbulence is Medium\\
          \textbf{AND} Action is Buy\\
          \textbf{THEN} Action is Hold

    \vspace{0.2cm}

    \item \textbf{IF} Turbulence is Medium\\
          \textbf{AND} Action is Sell\\
          \textbf{THEN} Action is Sell

    \vspace{0.2cm}

    \item \textbf{IF} Turbulence is Low\\
          \textbf{THEN} Action follows the original action
\end{enumerate}

\column{0.58\textwidth}

\begin{figure}
    \centering
    \includegraphics[width=\linewidth]
    {figures/action_smoother_mfs.png}
\end{figure}

\end{columns}

\end{frame}

% =====================================================
\begin{frame}{Fuzzy Reward Shaper}

\textbf{Goal}

Replace the fixed volatility penalty with a context-dependent fuzzy penalty.

\vspace{0.3cm}

Original reward:

\[
r_t=
\log\left(
\frac{V_{t+1}+\epsilon}
     {V_t+\epsilon}
\right)
-\lambda e_t v_t
\]

\vspace{0.3cm}

\textbf{Inputs}

\begin{itemize}
    \item Portfolio Exposure ($e_t$):
    \[
    \{\text{Low},\ \text{High}\}
    \]

    \item VIX Volatility ($v_t$):
    \[
    \{\text{Calm},\ \text{Stressed},\ \text{Panic}\}
    \]
\end{itemize}

\vspace{0.2cm}

\textbf{Output}

\[
\phi_t \in [0,1]
\]

Adaptive penalty factor.

\end{frame}

% =====================================================
\begin{frame}{Fuzzy Reward Shaper Rules}

\begin{columns}[T]

\column{0.42\textwidth}

\textbf{Rule Base}

\vspace{0.2cm}

\begin{enumerate}
    \item \textbf{IF} VIX is Panic\\
          \textbf{AND} Exposure is High\\
          \textbf{THEN} Penalty is Severe

    \vspace{0.3cm}

    \item \textbf{IF} VIX is Stressed\\
          \textbf{AND} Exposure is High\\
          \textbf{THEN} Penalty is Medium

    \vspace{0.3cm}

    \item \textbf{IF} VIX is Calm\\
          \textbf{OR} Exposure is Low\\
          \textbf{THEN} Penalty is None
\end{enumerate}

\vspace{0.3cm}

\[
r_t=
\log\left(
\frac{V_{t+1}+\epsilon}
     {V_t+\epsilon}
\right)
-
\phi_t\lambda_{scale}
\]

\column{0.58\textwidth}

\begin{figure}
    \centering
    \includegraphics[width=\linewidth]
    {figures/reward_shaper_mfs.png}
\end{figure}

\end{columns}

\end{frame}

% =====================================================
\begin{frame}{Genetic Algorithm Hyperparameter Optimization}

\textbf{Motivation}

\vspace{0.2cm}

DRL performance is highly sensitive to hyperparameter selection.

\vspace{0.2cm}

Examples include: Learning rate, Batch size, Rollout length etc

\vspace{0.3cm}

\textbf{Chromosome Representation}

Each hyperparameter configuration is encoded as a chromosome:

\[
C =
[g_1,g_2,\dots,g_n]
\]

where each gene corresponds to a specific hyperparameter.

\vspace{0.3cm}

\textbf{Fitness Function}

\[
Fitness = \text{Final Portfolio Value}
\]

obtained after training and evaluating the corresponding DRL agent.

\end{frame}

% =====================================================
\begin{frame}{Genetic Algorithm Evolution Cycle}

\textbf{Evolutionary Process}

\begin{enumerate}
    \item Generate an initial population of random chromosomes
    \item Train and evaluate each DRL agent
    \item Compute fitness using the final portfolio value
    \item Select the highest-performing individuals
    \item Apply crossover to generate offspring
    \item Apply mutation to maintain diversity
    \item Form the next generation and repeat
\end{enumerate}

\vspace{0.4cm}

\textbf{Genetic Operators}

\begin{itemize}
    \item \textbf{Selection}: favors high-fitness chromosomes
    \item \textbf{Elitism}: preserves the best solution unchanged
    \item \textbf{Uniform Crossover}: combines genes from two parents
    \item \textbf{Mutation}: randomly modifies genes to explore new regions
\end{itemize}

\vspace{0.2cm}

The cycle is repeated for multiple generations until the best
hyperparameter configuration is identified.

\end{frame}

% =====================================================
\begin{frame}[plain]

\begin{center}

\vspace{1cm}

{\Large \textbf{Thank You}}

\vspace{0.7cm}

{\large Questions?}

\vspace{1.2cm}

\textbf{Project Repository}

\vspace{0.2cm}

\url{https://github.com/georrous6/finrl-trading-lab}

\vfill

\small
Deep Reinforcement Learning Framework for Multi-Asset Trading

\end{center}

\end{frame}

\end{document}